\section{Multi-Flow Coexistence and Fairness Analysis}
\label{sec:fairness}

\subsection{Setup}
One subject flow shares one bottleneck with $k\in\{1,2,3\}$ background flows, so $n=2$--$4$ flows. The bottleneck is the UDP impairment proxy configured as $1000$ packets/s, a $60$-packet drop-tail queue, $30$\,ms one-way delay, $3$\,ms jitter and $0.1\%$ random loss; every flow sends $3000$ packets. The background flows are \textbf{user-space CUBIC and a simplified BBR-style controller running inside the same RUDP transport}, not kernel TCP. A flow's share is the number of its packets the bottleneck forwarded when the first flow finished. We report Jain's index $J=(\sum x_i)^2/(n\sum x_i^2)$ over all flows and the subject's share divided by its equal share $1/n$ (1.00 is fair). Every cell is $5$ runs with seeds $1$--$5$, the same seeds for every subject, so a standard deviation of $0.02$--$0.14$ on $J$ (\texttt{results/fairness/summary.csv}) is typical and differences below about $0.06$ should not be read as real.
Four subjects are compared: the static agent of Section~\ref{sec:mdp} (``Static''); the same agent with the inference-time \texttt{--polite} guard (``Static+guard''); the Poise agent of Section~\ref{sec:truesmart}; and an \emph{ablation} that is Poise trained identically but with the fairness weight $w_4=0$.
The Poise checkpoint was fixed from simulator validation before any real-network run, and all rows use the same code path, parameters and seeds.

\begin{table*}[t]
\centering\small
\setlength{\tabcolsep}{4pt}
\caption{Coexistence on the real UDP stack, mean of 5 runs [measured, proxy backend]. $J$: Jain index over all flows. subj: subject share / equal share. The subject is the policy in the column header.}
\label{tab:fair}
\begin{tabular}{llcc|cc|cc|cc}
\toprule
 & & \multicolumn{2}{c}{Static} & \multicolumn{2}{c}{Static+guard} & \multicolumn{2}{c}{No-fair.\ term} & \multicolumn{2}{c}{Poise} \\
bg & $k$ & $J$ & subj & $J$ & subj & $J$ & subj & $J$ & subj \\
\midrule
CUBIC & 1 & 0.957 & 1.20 & 0.911 & 0.72 & 0.879 & 1.37 & \textbf{0.982} & \textbf{0.88} \\
CUBIC & 2 & 0.853 & 1.58 & 0.940 & 0.80 & 0.788 & 1.72 & \textbf{0.969} & \textbf{1.09} \\
CUBIC & 3 & 0.827 & 1.73 & 0.951 & 0.75 & 0.708 & 2.08 & \textbf{0.968} & \textbf{1.03} \\
\midrule
BBR & 1 & 0.935 & 0.91 & 0.878 & 0.64 & 0.897 & 1.31 & \textbf{0.935} & \textbf{0.78} \\
BBR & 2 & 0.879 & 0.92 & 0.830 & 0.62 & 0.713 & 0.99 & \textbf{0.818} & \textbf{0.64} \\
BBR & 3 & 0.787 & 1.33 & 0.865 & 0.66 & 0.840 & 1.16 & \textbf{0.779} & \textbf{0.67} \\
\bottomrule
\end{tabular}
\end{table*}

\subsection{Results against CUBIC: the bully is gone}
Without a fairness term, the same network and training procedure produce a bully: the ablation takes $1.37\times$, $1.72\times$ and $2.08\times$ its share against one, two and three CUBIC flows, and $J$ falls to $0.708$. The static agent behaves similarly ($1.20\times$ to $1.73\times$). The Poise agent takes $0.88\times$, $1.09\times$ and $1.03\times$ its share and keeps $J$ at $0.982$, $0.969$ and $0.968$. Against the static agent the gain in $J$ is significant with two and three competitors ($0.969$ vs.\ $0.853$, $p=0.001$; $0.968$ vs.\ $0.827$, $p<0.001$, Welch) and so is the drop in the subject's share ($p\le0.001$ in all three cells). Against the fairness-ablated agent $J$ is higher at every $k$ ($p=0.012,\,0.002,\,0.014$): the fairness term, not the architecture, removes the bully behaviour. The smallest flow's share is $0.88$, $0.80$ and $0.79$ of fair for Poise, against $0.80$, $0.69$ and $0.57$ for the static agent.
Compared with the rule-based guard, Poise has a higher mean $J$ in all three CUBIC cells ($0.982$ vs.\ $0.911$, $0.969$ vs.\ $0.940$, $0.968$ vs.\ $0.951$) but these differences are not significant at five runs ($p=0.11,\,0.24,\,0.30$). The practical difference is the share: the guard over-corrects to $0.72$--$0.80\times$, whereas the learned policy sits at about $1\times$ (significant at $k=3$, $1.03$ vs.\ $0.75$, $p=0.018$). So the learned controller removes the need for the guard against CUBIC, and gives a bit more throughput than the guard does, but we cannot claim a measurable fairness gain over the guard from this sample size.

\subsection{Results against BBR-style flows: no improvement}
Against BBR-style competitors the result is negative. Poise has $J=0.935$, $0.818$ and $0.779$ for $k=1,2,3$, indistinguishable from the static agent ($0.935$, $0.879$, $0.787$; $p\ge0.24$) and from the guard ($0.878$, $0.830$, $0.865$; $p=0.051$ at $k=3$ in the guard's favour). It yields ($0.78\times$, $0.64\times$, $0.67\times$ of its share) and the smallest flow is only $0.44\times$ of fair at $k=3$. The same agent in the simulator reached $J\ge0.89$ against its simulated BBR flows, so the real BBR-style sender is harsher than our simulation of it. We tried retraining against randomised, more aggressive simulated BBR flows; in simulator validation the result was a milder bully (Jain $0.78$--$0.83$ against $0.88$) and we discarded it, so this gap remains open.

\subsection{Cost to a single flow}
Fairness must not be bought by slowing a lone flow. For a $256$\,kB transfer ($n=8$, byte-verified) Poise finished in $527$\,ms on a clean path against $599$\,ms for tuned AIMD ($p<0.001$) and $798$\,ms for the static agent ($p=0.053$); at $2\%$ loss in $1275$\,ms against $1572$ (AIMD, $p=0.052$) and $1263$ (static, $p=0.94$); at $10\%$ loss in $6004$\,ms against $6319$ (AIMD, $p=0.78$) and $7681$ (static, $p=0.34$). With FEC (one parity packet per eight data packets) it is $19\%$ slower on the clean path ($p=0.03$) and $17\%$ and $38\%$ faster at $2\%$ and $10\%$ loss ($p=0.047$, $0.045$); at $10\%$ loss the static agent with FEC is faster on average ($2398$ vs.\ $3743$\,ms, $p=0.11$). The lone-flow cost of Poise is therefore nil to negative on these paths, and the best single-flow result under heavy loss still belongs to the static agent with FEC.

\subsection{Limits of this evidence}
(i) The competitors are user-space controllers in RUDP, not kernel CUBIC or BBR. A kernel-TCP backend exists in \texttt{run\_fairness\_test.py} (netem on loopback, \texttt{TCP\_CONGESTION}) but was \textbf{not executed}: the development container has no \texttt{tc} or \texttt{ip}. (ii) Five runs per cell, one bottleneck configuration, equal flow sizes, and a measurement window set by the first flow to finish. (iii) Reward weights, hyper-parameters and the checkpoint were chosen by us on the simulator over about ten training runs; the real-network seeds ($1$--$5$) were not used for selection, but the BBR-randomisation experiment was prompted by seeing the BBR cells fail. (iv) The proxy never reorders packets or loses ACKs. (v) Slow start, pacing and the retransmission machinery are not learned.
